\chapter{Instrukcja instalacji i obsługi}

\section{Instalacja}

Repozytorium zawierające kod źródłowy znajduje się pod tym linkiem \cite{repo}. W celu uruchomienia aplikacji wymagane jest wcześniejsze zainstalowanie języka Python \cite{python}, a także dowolnego środowiska umożliwiającego edycję oraz wykonywanie programów napisanych w tym języku takiego jak na przykład PyCharm \cite{pycharm} lub Visual Studio Code \cite{visualstudiocode}. 

Następnie, w celu pobrania repozytorium na środowisko lokalne, należy otworzyć folder docelowy w terminalu i użyć następującej komendy:

\hfill \break
\textbf{git clone https://github.com/Sandpitturtleee/Thesis.git}\newline
\hfill 

Oprócz tego konieczne jest również zainstalowanie wszystkich wymaganych zewnętrznych paczek. Lista wykorzystywanych bibliotek wraz z ich wersjami znajduje się w pliku \textbf{requirements.txt}. Ich instalację można przeprowadzić za pomocą następującej komendy:

\hfill \break
\textbf{pip install -r requirements.txt}\newline
\hfill 

\section{Instrukcja obsługi}

Projekt został podzielony na trzy główne części odpowiadające kolejnym jego etapom. Jest to generowanie grafów, wykonywanie algorytmów Dijkstry oraz analiza otrzymanych wyników. Parametry wykorzystywane podczas działania programu znajdują się w pliku \textbf{config.py}.

\subsection{Generowanie grafów}

Pierwszym etapem działania aplikacji jest przygotowanie grafów wykorzystywanych następnie podczas testowania algorytmów. Kod odpowiedzialny za ten proces znajduje się w katalogu \path{graphs_creation/src}. Głównym plikiem służącym do uruchomienia procesu generowania grafów jest \path{graphs_creation/src/main.py}. Program generuje grafy o parametrach określonych w konfiguracji projektu. W pliku \path{config.py} możliwe jest określenie maksymalnego rozmiaru grafu oraz liczby generowanych przypadków. 

\subsection{Analiza grafów}

Kolejnym etapem jest wykonanie algorytmów wyznaczania najkrótszych ścieżek Dijkstry dla wcześniej przygotowanych grafów. Kod odpowiedzialny za ten proces znajduje się w katalogu \path{graphs_analysis/src}. Główny plik uruchamiający analizę znajduje się pod ścieżką \path{graphs_analysis/src/main.py}. W katalogu \path{graphs_analysis/src} implementacje zostały rozdzielone na dwie części: \path{standard} oraz \path{quantum}. Pierwsza z nich zawiera klasyczne implementacje algorytmu Dijkstry, natomiast druga odpowiada za implementację kwantową.

\subsection{Analiza statystyczna wyników}

Po przeprowadzeniu eksperymentów otrzymane wyniki mogą zostać poddane dalszej analizie. Do jej przeprowadzenia wykorzystywany jest plik \path{data_analysis/src/main_analysis.py} znajdujący się w katalogu \path{data_analysis/src}.

\subsection{Generowanie wykresów}

W celu graficznego przedstawienia rezultatów przeprowadzonych eksperymentów w projekcie przygotowany został osobny moduł odpowiedzialny za tworzenie wykresów. Głównym plikiem wykorzystywanym w tym celu jest \path{data_analysis/src/main_plots.py}. Po wykonaniu generowane są wykresy pozwalające na porównanie rezultatów otrzymanych dla różnych implementacji algorytmu oraz różnych rodzajów i rozmiarów grafów. 

